% !TeX root = ../navier-stokes-manuscript.tex
% ============================================================
% 3.1 The Terminal Quantifier
% ============================================================
\subsection{The Terminal Quantifier}\label{sub:TheTerminalQuantifier}

Keep the rapid-decay datum \(u_0\), the same viscosity \(\nu>0\), and the maximal classical velocity and pressure history \((u,p)\) that they generate.
Let its maximal lifespan be \([0,T_*)\), suppose for contradiction that \(T_*<\infty\), and write \(I_*=(0,T_*)\).
Every derivative introduced below is a coordinate of this one history.
The datum is smooth and divergence-free, so it belongs to every finite solenoidal Sobolev space.
Write
\begin{equation}\label{eq:SectionThreeSmoothSolenoidalClass}
  \Hinfty_\sigma(\R^3)
  :=
  \left\{
    v\in\bigcap_{m\in\Nzero}H^m(\R^3;\R^3)
    :\nabla\cdot v=0
  \right\}.
\end{equation}

Because the history decays at spatial infinity, the whole-space normalization fixes pressure from the same velocity field.
If \(R_j=\partial_j(-\Delta)^{-1/2}\) denotes the \(j\)-th Riesz transform, then
\begin{equation}\label{eq:SectionThreePressureNormalization}
  -\Delta p=\partial_i\partial_j(u_i u_j),
  \qquad
  p=R_iR_j(u_i u_j).
\end{equation}
For \(V_n=\partial_t^n u\) and \(Q_n=\partial_t^n p\), direct differentiation of the original equation gives
\begin{align}
  Q_n
  &=R_iR_j\sum_{a=0}^{n}\binom{n}{a}V_{a,i}V_{n-a,j},
  \label{eq:SectionThreePressureRecurrence}\\
  V_{n+1}
  &=\nu\Delta V_n
  -\sum_{a=0}^{n}\binom{n}{a}(V_a\cdot\nabla)V_{n-a}
  -\nabla Q_n.
  \label{eq:SectionThreeVelocityRecurrence}
\end{align}
The binomial sum retains every way that the time derivatives divide across the transported and transporting velocity factors.
At every temporal depth, transport, pressure, viscosity, and incompressibility remain joined inside the same evolution.

To inspect a finite part of the tower, choose independent temporal and spatial depths \(N\) and \(M\).
For finite \(N,M\in\Nzero\) and \(0<T<T_*\), define the adjacent-row rectangle norm
\begin{equation}\label{eq:SectionThreeRectangleNorm}
\begin{aligned}
  \mathfrak R_{N,M}(u;T)
  :={}&
  \sum_{n=0}^{N}
  \norm{\partial_t^n u}_{L^2(0,T;H^{M+1})}^{2}
  \\
  &+
  \sum_{n=0}^{N}
  \norm{\partial_t^{n+1}u}_{L^2(0,T;H^{M-1})}^{2}.
\end{aligned}
\end{equation}
Negative Sobolev indices are understood in the inhomogeneous Bessel-potential scale.
The finite record \(\cR_{N,M}[u_0;T]\) consists of these velocity rows together with the pressure rows and differentiated equations that identify them as derivatives of the same history.
Classical smoothness already gives every such record after a strict terminal time \(T<T_*\) has been fixed.

\begin{lemma}[Rectangles on strict classical subintervals]\label{lem:StrictSubintervalRectangles}
For every \(T<T_*\) and every finite \(N,M\), one has
\begin{equation}\label{eq:AppendixStrictSubintervalRectangleBound}
  \mathfrak R_{N,M}(u;T)<\infty.
\end{equation}
If \(N'\leq N\), \(M'\leq M\), and \(T'\leq T\), restriction of the larger rectangle gives the smaller rectangle.
\end{lemma}

\begin{proof}
The closed interval \([0,T]\) lies inside the classical lifespan.
Smooth-data persistence gives \(\partial_t^n u\in C([0,T];H^q(\R^3))\) for every finite temporal order \(n\) and every finite Sobolev order \(q\).
Compactness of \([0,T]\) makes each of these norms bounded.
In particular,
\[
  \norm{\partial_t^n u}_{L^2(0,T;H^q)}^2
  \leq
  T\sup_{0\leq t\leq T}
  \norm{\partial_t^n u(t)}_{H^q}^{2}
  <\infty.
\]
Choosing the spatial orders in \eqref{eq:SectionThreeRectangleNorm} proves \eqref{eq:AppendixStrictSubintervalRectangleBound}.
Restriction preserves the solution and every one of its derivatives, so every shared coordinate agrees.
\end{proof}

The strict-subinterval lemma stops at \(T<T_*\), and its bound may grow as \(T\uparrow T_*\).
The whole-terminal estimate in Theorem~\ref{thm:BodyWholeTerminalRectangleTheorem}
supplies a bound independent of the cutoff after \(N\) and \(M\) have been fixed.
To state the agreement between different finite depths, define, for \(N'\leq N\) and \(M'\leq M\), the terminal bonding map
\[
  \rho_{N',M';I_*}^{N,M;I_*}:
  \mathscr X_{N,M}(I_*)
  \longrightarrow
  \mathscr X_{N',M'}(I_*)
\]
by forgetting the higher rows, applying the Sobolev embeddings to every retained velocity and pressure row, and recomputing the scalar adjacent-row norm on \(I_*\).
Forgetting in two stages gives the same result as forgetting directly, and leaving the depths unchanged gives the identity map.

\begin{theorem}[Datum-generated whole-terminal estimate in rectangle form]\label{thm:WholeTerminalCompatibleRectangleEstimate}
For every finite \(N,M\in\Nzero\), the same datum-generated pressure-complete history satisfies
\begin{equation}\label{eq:WholeTerminalRectangleBound}
\begin{aligned}
  \mathfrak R_{N,M}(u;I_*)
  :={}&
  \sum_{n=0}^{N}
  \norm{\partial_t^n u}_{L^2(I_*;H^{M+1})}^{2}
  \\
  &+
  \sum_{n=0}^{N}
  \norm{\partial_t^{n+1}u}_{L^2(I_*;H^{M-1})}^{2}
  <\infty.
\end{aligned}
\end{equation}
This finite norm supplies the whole-terminal record
\[
  \cR_{N,M}[u_0]
  \in
  \mathscr X_{N,M}(I_*),
\]
obtained from the same velocity and normalized pressure rows on \(I_*\).
If \(N'\leq N\) and \(M'\leq M\), the terminal restriction map forgets the higher temporal coordinates, applies the continuous Sobolev embeddings, and satisfies
\begin{equation}\label{eq:WholeTerminalRectangleCompatibility}
  \rho_{N',M';I_*}^{N,M;I_*}
  \bigl(\cR_{N,M}[u_0]\bigr)
  =\cR_{N',M'}[u_0].
\end{equation}
Every shared coordinate is the corresponding distributional derivative of the same \(u\), and every pressure row uses the normalization in \eqref{eq:SectionThreePressureNormalization}.
\end{theorem}

\begin{remark}[Coordinate realization of the whole-terminal estimate]
Fix \(N,M\).
For every \(0\leq n\leq N\), Theorem~\ref{thm:BodyWholeTerminalRectangleTheorem} gives \(V_n\in L^2(I_*;H^{M+1})\) and \(V_{n+1}\in L^2(I_*;H^{M-1})\), all from the same pressure-complete history.
Adding those finitely many squared norms gives the velocity part of \eqref{eq:WholeTerminalRectangleBound}.
Riesz-transform boundedness, the Sobolev product estimate, and Cauchy--Schwarz in time give, for \(0\leq n\leq N\),
\[
  \norm{Q_n}_{L^1(I_*;H^M)}
  \leq
  C_{M,n}\sum_{a=0}^{n}
  \norm{V_a}_{L^2(I_*;H^{M+1})}
  \norm{V_{n-a}}_{L^2(I_*;H^{M+1})}
  <\infty.
\]
The pressure estimate completes \eqref{eq:WholeTerminalRectangleBound}.
Forgetting higher temporal rows or lowering the spatial index leaves every shared distributional derivative and pressure row unchanged, which gives \eqref{eq:WholeTerminalRectangleCompatibility}.
\end{remark}

Conversely, the finite-rectangle form recovers the adjacent rows coordinate by coordinate.
\begin{proposition}[Adjacent-row recovery]\label{prop:AdjacentRowRecovery}
The whole-terminal estimate gives, for every finite \(n,m\in\Nzero\),
\begin{equation}\label{eq:WholeTerminalAdjacentRows}
  V_n\in L^2(I_*;H^{m+1}),
  \qquad
  \partial_tV_n=V_{n+1}\in L^2(I_*;H^{m-1}).
\end{equation}
Each finite pair has its own constant.
\end{proposition}

\begin{proof}
Fix \(n\) and \(m\), and choose the rectangle with \(N=n\) and \(M=m\).
The two terms indexed by \(n\) in the whole-terminal sum \eqref{eq:WholeTerminalRectangleBound} are finite, which gives the memberships in \eqref{eq:WholeTerminalAdjacentRows}.
Every row in \eqref{eq:WholeTerminalAdjacentRows} is an actual derivative of \(u\), and \eqref{eq:WholeTerminalRectangleCompatibility} preserves that identity under restriction.
\end{proof}

The strict and whole-terminal rectangles use the same distributional derivatives of \((u,p)\).
Equation \eqref{eq:WholeTerminalRectangleCompatibility} makes the full family compatible under every decrease of \(N\) or \(M\).
